\section*{Chapter \ref{ch:b3:hydrogen-atom} --- The Hydrogen Atom}

\begin{solution}{exo:b3:hydrogen-atom:1}
(a) $u'' = (1/a^2 - 2/ar)u$; the $1/r$ terms match when
$\hbar^2/m_{\text{e}}a = k$, i.e.\ $a = a_0$, and the constant terms
give $E = -\hbar^2/2m_{\text{e}}a_0^2 = -E_{\text{I}}$. (b) $C =
2/a_0^{3/2}$. (c) $[\hbar^2/mk] = \unit{m}$; $[mk^2/\hbar^2] =
\unit{J}$. (d) The variational argument of
\cref{exo:b3:schrodinger-three-dimensions:12} shows $-E_{\text{I}}$
is the least energy any normalised state can achieve: the
uncertainty principle floors the atom.
\end{solution}

\begin{solution}{exo:b3:hydrogen-atom:2}
(a) \qty{121.6}{nm}; limit \qty{91.2}{nm}. (b) $656.3$ (red),
$486.1$ (blue-green), \qty{434.0}{nm} (violet); limit
\qty{364.6}{nm} --- the first three are visible. (c) From
\qty{1875}{nm} down to \qty{820}{nm}: near infrared. (d)
$R_{\text{H}}c\,(1/109^2 - 1/110^2) = \qty{5.01}{GHz}$.
\end{solution}

\begin{solution}{exo:b3:hydrogen-atom:3}
(a) $\dd(r^2\eu^{-2r/a_0})/\dd r = 0$ at $r = a_0$. (b) $\langle r\rangle = (4/a_0^3)\int_0^\infty
r^3\eu^{-2r/a_0}\dd r = (4/a_0^3)\times 3!\,(a_0/2)^4 = \tfrac32
a_0$. (c) With $x = 4$:
$\tfrac12(16 + 8 + 2)\eu^{-4} = 0.24$: a quarter of the time the
electron is beyond $2a_0$. (d) The electron has no radius: $P(r)$
is a distribution with a mode, a mean and long tails --- the atom
is fuzzy at the factor-two level.
\end{solution}

\begin{solution}{exo:b3:hydrogen-atom:4}
(a) $l = 0$: one; $l = 1$: three; $l = 2$: five --- nine. (b) $2$,
$8$, $18$. (c) Exactly the lengths of the first three rows: the
periodic table is the filling record of these shells. (d) The
$m$-degeneracy survives (space is still isotropic); the
$l$-degeneracy breaks because the screened potential seen by the
valence electron is no longer pure $1/r$ ---
\cref{exo:b3:hydrogen-atom:8}.
\end{solution}

\begin{solution}{exo:b3:hydrogen-atom:5}
(a) Hydrogen-like with $Z - 1$: $h\nu = (Z-1)^2E_{\text{I}}(1 -
\tfrac14)$. (b) $784 \times \qty{10.2}{eV} = \qty{8.0}{keV}$ ---
matching the measured K$\alpha$ of copper. (c) That the integer
ordering the elements is the \emph{nuclear charge}, not the atomic
weight: gaps in Moseley's lines located undiscovered elements. (d)
$42^2 \times 10.2 = \qty{18.0}{keV}$ --- technetium, found decades
later, obliged.
\end{solution}

\begin{solution}{exo:b3:hydrogen-atom:6}
(a) Two equal masses: $\mu = m_{\text{e}}/2$: binding
\qty{6.8}{eV}, radius $2a_0 = \qty{0.106}{nm}$. (b) $\tfrac34
\times 6.8 = \qty{5.1}{eV}$: \qty{243}{nm}. (c) Two photons of
\qty{511}{keV}, back to back (momentum conservation at rest). (d)
Positron-emission tomography: the pair of collinear
\qty{511}{keV} photons is the signal every PET ring triangulates.
\end{solution}

\begin{solution}{exo:b3:hydrogen-atom:7}
(a) $a_{\text{ex}} = 0.0529 \times 12.9/0.058 = \qty{11.8}{nm}$.
(b) $E = 13.6 \times 0.058/12.9^2 = \qty{4.7}{meV}$. (c) The
natural pair size exceeds the dot: the wall, not the attraction,
shapes the state --- the $1/R^2$ beats the $1/R$, as promised. (d)
$k_{\text{B}}T$ at room temperature is five times the binding: the
pairs ionise; excitonic physics lives below $\sim\qty{50}{K}$ in
clean crystals.
\end{solution}

\begin{solution}{exo:b3:hydrogen-atom:8}
(a) $3s$: its innermost lobe dives inside the core where the
nuclear $+11$ is barely screened. (b) $-13.6/9 = \qty{-1.51}{eV}$
matches $3d$ almost exactly: $3d$ orbits outside the core and sees
a net $+1$; $3s$ (and partly $3p$) taste the deep potential and
sink. (c) $\Delta E = \qty{2.10}{eV}$: $\lambda = \qty{590}{nm}$
--- the sodium yellow of street lamps. (d) $n_{\text{eff}} =
\sqrt{13.6/|E|}$: $1.63$ and $2.11$, so $\delta_s = 1.37$,
$\delta_p = 0.89$.
\end{solution}

\begin{solution}{exo:b3:hydrogen-atom:9}
(a) About the binding of $n = 100$, $\sim\qty{1.4}{meV}$, plus the
electron's small thermal energy: a far-infrared/radio photon. (b)
The $\Delta n = 1$ steps near $n \approx 100$ fall in the GHz radio
band --- the recombination lines. (c) H$\alpha$ is the step $3 \to
2$, near the cascade's end. (d) Lyman $\alpha$, though strongest,
is ultraviolet and, in a nebula full of ground-state hydrogen, is
resonantly scattered and trapped; H$\alpha$ escapes freely and
paints the nebula red.
\end{solution}

\begin{solution}{exo:b3:hydrogen-atom:10}
(a) $\langle 1/r\rangle = 1/a_0$: $\langle E_p\rangle = -k/a_0 =
\qty{-27.2}{eV} = 2E_1$. (b) $\langle E_k\rangle = E_1 - \langle
E_p\rangle = +\qty{13.6}{eV}$: the virial ratio $-\tfrac12$ of
every $1/r$ orbit. (c) $v = \sqrt{2E_k/m_{\text{e}}} = \alpha c =
\qty{2.2e6}{m/s}$. (d) $I = ev/2\pi a_0 \approx \qty{1}{mA}$
circulating at \qty{53}{pm}: $B \sim \mu_0I/2a_0 \approx
\qty{12}{T}$ --- the enormous internal field hyperfine structure
will feed on.
\end{solution}

\begin{solution}{exo:b3:hydrogen-atom:11}
(a) From \cref{exo:b3:hydrogen-atom:10}(c), $v/c = \alpha$;
scaling: $Z\alpha/n$. (b) $\alpha^2E_{\text{I}} = \qty{0.72}{meV}$;
at $n = 2$ the splittings come out tens of $\unit{\micro eV}$ ---
some \qty{10}{GHz}, radio-measurable (and measured). (c) $Z\alpha =
0.67$: two-thirds of light speed --- the Schrödinger treatment
fails; Dirac's equation takes over. (d) At $Z\alpha \to 1$ the
ground state's energy formally dives past $-2m_{\text{e}}c^2$: the
vacuum itself would spark electron--positron pairs --- supercritical
fields, sought in heavy-ion collisions.
\end{solution}

\begin{solution}{exo:b3:hydrogen-atom:12}
(a) Differentiate: $\dot{\vect p}\wedge\vect L = -(mk/r^3)\vect
r\wedge(\vect r\wedge\dot{\vect r}\,m)$; expanding the double cross
product gives exactly $mk\,\dd\vect e_r/\dd t$: $\dot{\vect A} =
\vect 0$. (b) A conserved axis is a symmetry beyond isotropy;
degeneracy is incomplete labelling, and the extra label ($\vect
A$'s quantum cousin) connects the $l$'s within one $n$. (c) With
$V = -k/r + \epsilon/r^2$ the effective angular momentum shifts,
the angular period no longer matches the radial one, and the
ellipse's axis turns at a rate $\propto\epsilon$. (d) Relativity's
corrections play the role of $\epsilon$ in hydrogen itself: the
fine structure of \cref{exo:b3:hydrogen-atom:11} is precisely the
$l$-degeneracy breaking.
\end{solution}

\begin{solution}{pb:b3:hydrogen-atom:1}
\textbf{1.} $\langle r\rangle \sim n^2a_0$; $|E_n| =
E_{\text{I}}/n^2$; spacing $\approx 2E_{\text{I}}/n^3$; orbital
frequency $\propto 1/n^3$ (Kepler on the Bohr orbit).
\textbf{2.} $E_{n+1} - E_n \to 2E_{\text{I}}/n^3$, and $h$ times
the classical frequency $v_n/2\pi r_n$ equals the same expression
(the computation of \cref{pb:b3:hamiltonian-mechanics:1}, item 20).
\textbf{3.} $n = 50$: $\qty{132}{nm}$, \qty{5.4}{meV}; $n = 100$:
$\qty{0.53}{\micro m}$, \qty{1.4}{meV}.
\textbf{4.} $n^2ea_0 = \qty{2.1e-26}{C.m} \approx 6300\,\text{D}$:
three thousand water molecules' worth of dipole on one atom.
\textbf{5.} $F \sim \num{5.4e-3}/\num{1.32e-7} \approx
\qty{4e4}{V/m} = \qty{400}{V/cm}$; the exact classical-ionisation
threshold is about eight times smaller ($\sim\qty{50}{V/cm}$):
either way, a whisper of a field.
\textbf{6.} With geometric cross-sections $\sim\pi n^4a_0^2$, even
ultra-high laboratory vacuum collides such an atom in
microseconds; interstellar densities ($\sim\num{e6}$ particles per
$\unit{m^3}$) leave it days --- space is the better vacuum chamber.
\textbf{7.} $\nu = R_{\text{H}}c\,(1/109^2 - 1/110^2) =
\qty{5.01}{GHz}$.
\textbf{8.} Centimetre radio: a radio telescope --- a big dish and
a quiet receiver.
\textbf{9.} $\lambda = \qty{6.0}{cm}$; the $n = 110$ atom is
$\sim\qty{0.6}{\micro m}$: a hundred thousand atoms per
wavelength --- comfortably an ``antenna'' regime.
\textbf{10.} $\Delta\nu/\nu \sim v/c = \num{4.3e-5}$: about
\qty{200}{kHz} on \qty{5}{GHz} --- and the measured width hands
back the nebular temperature.
\textbf{11.} The ionised inner Galaxy: HII regions hidden behind
dust that extinguishes every Balmer photon --- radio recombination
lines mapped the spiral structure optical astronomy could not see.
\textbf{12.} The micro electric fields of neighbouring ions (Stark
broadening) smear the levels, and blackbody/cosmic radiation plus
collisions ionise or $l$-mix the fragile states.
\textbf{13.} $C_6 \sim d^4/\Delta E \propto (n^2)^4/n^{-3} =
n^{11}$: raise $n$ from $1$ to $70$ and the interaction grows by
twenty orders of magnitude.
\textbf{14.} Within $R_{\text{b}}$ the pair state is shifted off
resonance, so one atom's excitation \emph{conditions} its
neighbour's response: exactly the controlled logic a two-qubit
gate needs.
\textbf{15.} Micrometre-range interactions let each atom sit in
its own addressable tweezer, far apart by atomic standards yet
strongly coupled --- interaction range matched to optical
resolution.
\textbf{16.} $\sim 200$ gates per lifetime: the ratio bounds the
achievable fidelity, and it is the ratio, not the raw
microseconds, that a processor lives on.
\textbf{17.} At \qty{300}{K} the photon bath is dense at
$\sim\qty{0.1}{meV}$: it shuffles and ionises Rydberg levels within
their radiative lifetimes; cold (\qty{4}{K}) shields starve those
transitions.
\textbf{18.} Transition rates scale as $d^2 \propto n^4$: at $n =
50$, six million times an ordinary atom's coupling --- which is why
a vapour cell of Rydberg atoms is now a calibrated microwave and
terahertz field sensor.
\textbf{19.} $92^2 \times \qty{13.6}{eV} = \qty{115}{keV}$ down to
$13.6/10^6\,\unit{eV} = \qty{13.6}{\micro eV}$: ten orders of
magnitude from one formula.
\textbf{20.} Classicality dawns where the spacing becomes a
vanishing fraction of the energy, $n \gg 1$ --- the atom's own
correspondence limit.
\textbf{21.} It is the one atom computable from first principles
to the experiment's precision: any mismatch is discovery, so
hydrogen anchors the fundamental constants.
\textbf{22.} $\tau \approx 1.6\,\unit{ns} \times 50^3 =
\qty{0.2}{ms}$ --- generous, but blackbody redistribution (Part
III) eats into it, another reason for the cold shields.
\textbf{23.} CPT symmetry --- matter and antimatter atoms must
match line by line; hydrogen is the anchor because only there does
theory reach the fifteenth digit alongside experiment.
\textbf{24.} The $n^2$ dipole makes the atom feel a single
microwave photon's field, while its long lifetime and level
structure let it acquire a measurable phase \emph{without}
absorbing the photon: quantum non-demolition sensing.
\textbf{25.} Scalings $n^2$, $n^{-2}$, $n^{-3}$, $n^{11}$ stretch
the one solved atom from \qty{52.9}{pm} to half a micrometre and
from \qty{121.6}{nm} to \qty{6}{cm} --- the same Schrödinger
solution broadcasting from the Galaxy's HII regions and clocking
two-qubit gates in tweezer arrays.
\end{solution}
